% Generated draft: requires mathematical and source review.
\documentclass[11pt]{article}
\usepackage[margin=25mm]{geometry}
\usepackage{fontspec}
\setmainfont{Latin Modern Roman}
\usepackage{amsmath,amssymb,mathtools,unicode-math}
\setmathfont{Latin Modern Math}
\usepackage{enumitem,needspace,xurl,hyperref}
\setlength{\parindent}{0pt}
\setlength{\parskip}{5pt}
\setlength{\emergencystretch}{3em}
\newcommand{\N}{\mathbb{N}}
\newcommand{\Z}{\mathbb{Z}}
\newcommand{\Q}{\mathbb{Q}}
\newcommand{\R}{\mathbb{R}}
\newcommand{\C}{\mathbb{C}}
\newcommand{\F}{\mathbb{F}}
\newcommand{\A}{\mathbb{A}}
\newcommand{\PP}{\mathbb{P}}
\newcommand{\E}{\mathbb{E}}
\newcommand{\eps}{\varepsilon}
\newcommand{\dd}{\,\mathrm{d}}
\begin{document}
\section*{UnsolvedMath ID: 701428}
\textbf{Source problem code:} catalogue entry
\subsection{Definitions and mathematical statement}
Let \(p\) be a prime. A finite group \(G\) is soluble of derived length \(d\) if its derived series, defined by \(G^{(0)}=G\) and \(G^{(i+1)}=[G^{(i)},G^{(i)}]\), first reaches the trivial group at \(G^{(d)}\). An elementary abelian \(p\)-group of order \(p^n\) is a group isomorphic to \((\Z/p\Z)^n\). For a subgroup \(A\leq\operatorname{Aut}(G)\), define its common fixed-point subgroup by
\[
C_G(A)=\{g\in G: a(g)=g\text{ for every }a\in A\}.
\]
For a group \(H\), put \(\gamma_1(H)=H\) and \(\gamma_{j+1}(H)=[\gamma_j(H),H]\). Saying that \(H\) is nilpotent of class at most \(c\) means that \(\gamma_{c+1}(H)=1\); the trivial group has class zero.

The question is whether there exist functions
\[
f:\Z_{\geq0}\longrightarrow\Z_{\geq0},\qquad
h:\{\text{primes}\}\times\Z_{\geq0}\times\Z_{\geq0}\longrightarrow\Z_{\geq0}
\]
such that, for every prime \(p\), every finite soluble group \(G\) of derived length \(d\), and every elementary abelian subgroup \(A\leq\operatorname{Aut}(G)\) of order \(p^n\) satisfying \(C_G(A)=1\), there is a normal series
\[
1=G_0\trianglelefteq G_1\trianglelefteq\cdots\trianglelefteq G_k=G,
\qquad k\leq f(n),
\]
whose factors satisfy
\[
\gamma_{h(p,n,d)+1}(G_i/G_{i-1})=1
\qquad(1\leq i\leq k).
\]
Thus the length bound must depend only on \(n\), and the bound on the nilpotency classes must depend only on \(p,n,d\), uniformly over \(G\) and \(A\). No assumption that \(p\) is coprime to \(|G|\) is supplied. The cases \(n=0\) or \(d=0\) force \(G=1\), for which the series of length zero is allowed.

There is a convention requiring clarification: the displayed chain expresses normality of each term in the next. If the source uses ``normal series'' to require every term to be normal in \(G\), the additional requirement is \(G_i\trianglelefteq G\) for all \(i\). The supplied record does not explicitly distinguish these conventions.

The supplied source states that the answer is affirmative for \(p=2\), and that an affirmative answer would follow from an affirmative answer to Kourovka Notebook Problem~11.125. These are retained as source assertions; the formulation of Problem~11.125 and the supporting arguments are not supplied.
\subsection{Short English statement}
Does a finite soluble group with a fixed-point-free elementary abelian group of automorphisms of order \(p^n\) have a normal series whose length is bounded solely by \(n\), with nilpotent factors of class bounded solely by \(p,n,d\), where \(d\) is its derived length? Here fixed-point-free means that the common fixed-point subgroup is trivial. The supplied source asserts the result for \(p=2\); the convention for normal series requires clarification.
\subsection{Sources}
\begin{itemize}
\item Supplied local UnsolvedMath registry record, permanent ID \(701428\), exact source code \texttt{catalogue entry}, titled ``Kourovka Notebook Problem 17.127.'' The supplied status is \texttt{open}; this is a registry label, not an independently verified assessment. Upstream revision: \texttt{d371ee917ac9647fbb06cd6fc529deec6f3dda82}. Record updated \texttt{2026-10-16T14:00:00Z}.
\item \emph{The Kourovka Notebook (Unsolved Problems in Group Theory)}, Problem~17.127, proposed by P. Shumyatsky in the 17th issue (2010). The supplied background also identifies the 21st issue edition (2026), \texttt{arXiv:1401.0300v48}. The registry collection attributes editorship to E. I. Khukhro and V. D. Mazurov. These bibliographic details are supplied attribution and have not been independently checked.
\item The supplied problem text states that the assertion holds for \(p=2\) and refers to Kourovka Notebook Problem~11.125 as a sufficient affirmative antecedent. No separate supporting reference or text of that problem is supplied.
\end{itemize}
\end{document}
